
\section{Background}
\label{sec:background}

\subsection{Tool-Enabled LLM Agents and Metadata}

Tool-enabled LLM agents augment a language model with external functionality,
such as filesystem access, Web retrieval, database queries, and code execution.
A typical system consists of an \emph{LLM}, an \emph{orchestration layer}, and
a set of \emph{tools}. The orchestrator exposes each tool to the LLM through
metadata describing its purpose, arguments, primitive types, and
natural-language semantics, then translates model outputs into structured
invocations and dispatches them to the corresponding implementations.
Importantly, schemas generally capture only syntactic constraints, whereas
tool descriptions may specify stronger semantic requirements. For example, a
\texttt{path} argument may be typed as a string but described as a ``file
within the workspace,'' implying a filesystem-containment property absent from
the schema. Similarly, strings representing URLs, commands, queries, or
hostnames require different security constraints despite sharing the same
primitive type. Thus, tool metadata provides a semantic specification of
LLM-generated arguments that ultimately become ordinary program inputs at the
tool boundary.

\subsection{Fuzzing Semantic Security Bugs}

Traditional fuzzers detect bugs using observable failures such as crashes,
assertion violations, sanitizer reports, or abnormal termination. Tool
implementations can instead execute successfully while violating their
documented semantics. For example, a workspace-local file tool may accept
\texttt{../../etc/passwd}, successfully access the file, and return normally.
We call such failures \emph{semantic security bugs}: executions that satisfy a
tool's syntactic interface but violate a security property implied by its
documented semantics.

Detecting these bugs requires semantic oracles rather than crash-based ones.
The required property depends on how an input is consumed: paths require
containment and traversal checks, URLs may require restrictions on schemes and
destinations, and command arguments require checks on interpretation by the
shell. Consequently, there is no universal sanitizer even for inputs sharing
the same primitive type. CWE captures recurring classes of such weaknesses,
including path traversal, command injection, improper input validation, and
server-side request forgery, but determining which checks apply normally
requires program-specific knowledge. Tool-enabled LLM systems expose an
unusual opportunity: the argument descriptions already supplied to the LLM
encode much of this semantic information, allowing security checks and
corresponding fuzzing oracles to be specialized to individual tool arguments.


\section{Threat Model}
\label{sec:threat-model}

\project{} assumes that the orchestration layer faithfully exposes a tool's
declared interface and dispatches schema-valid invocations to the corresponding
implementation. We trust the \project{} analysis infrastructure, but do not
trust the LLM or the tool implementation. In particular, tool metadata
describes the behavior expected by the LLM, while the implementation may fail
to enforce the corresponding security properties.

We consider an attacker who can influence information consumed by the LLM,
including user prompts, retrieved content, and outputs of other tools. Because
such information may propagate into structured tool arguments, we treat every
schema-valid argument as potentially adversarial. A vulnerability exists when
the tool accepts or processes such an argument in a way that violates a
security property implied by its documented interface. Vulnerability discovery
is independent of a particular LLM or prompting strategy; we separately
evaluate whether discovered triggering inputs can be delivered through the
normal LLM tool-calling path.

